\cotefile{93}

\scanpage{76}

\unsure{b'|b)|Si} on s'intéresse à $\unsure{\mathscr E|\mathscr L|\mathcal E}_X$ au lieu de $\underline O_X$, il faut remplacer $\underline\omega_X$ par $\underline\omega_X^*$ (sous-faisceau des bases de $\underline\omega_X$, i.e. des sections sans zéros), qui est un torseur sous $\underline O_X^*$ et $dg-df$ par $\deleted{dg}\mathrel{:} A'(f,g)=\frac{dg}{df}\in\Gamma(X,\underline O_X^{\times})$

b) \uline{$G_m$ opérant sur $\mathbb E^{\unsure{1|\prime}}$ ou $\mathbb E^{1\times}\simeq G_m$} par \deleted{translation}\inserted{homothétie}. Plus généralement, considérons son opération sur un faisceau inversible $L$, soit $L^{\times}$ le sous faisceau des bases \inserted{(torseur sous $\underline O_X^{\times}$)}, on s'intéresse à $L^{\times}/\mathbb C^{\times}$. Or

à \unsure{tte|une} section de $L^{\times}$ est \unsure{associé|associé'|associée} une connexion sur $L$. on trouve ainsi un \unsure{hom.|hom|homom.}

\[
 L^{\times}\,\deleted{\unsure{/}}\xrightarrow{D}\underline{\mathrm{Conn}}_{X/\mathbb C}(L)\qquad(f\longmapsto D_f)
\]

où le deuxième mb est un torseur sous $\underline\omega_{X/\mathbb C}^{\unsure{1}}$, et ceci induit une \deleted{\unsure{bijection}} Iso

\[
 L^{\times}/\mathbb C^{\times}\simeq\underline{\mathrm{Conn}}_{X/\mathbb C}(L)
\]

Si on ignorait la notion de connexion, on introduirait directement $B(f,g)$ $(f,g\in\Gamma L^*)$, qui sera \deleted{\unsure{comme fonct.|conn.}} $D_g-D_f$, par

\[
 B(f,g)=\frac{d\,g/f}{g/f}\in\Gamma\deleted{(}\underline\omega_{X/\mathbb C}
\]

avec des propriétés du type a) b) c) d) e). Dans le cas $L=\underline O_X$, $\underline{\mathrm{Conn}}(L)\simeq\underline\omega^{\unsure{1}}_{X/\mathbb C}$ can$^{\mathrm{t}}$, et il suffit de considérer

\[
 g\longmapsto\frac{dg}{g}\qquad O_X^*\longrightarrow\underline\omega_{X/\mathbb C}
\]

qui induit $\underline O_{X}^{\times}/\mathbb C^{\times}\simeq\underline\omega_{X/\mathbb C}$

% archive sheet number: 75
